% Lecture 1 notes: \input by notes.tex, \subimport'ed by the semester book.
% Exam questions (зачёт): none (introductory lecture; DLP/DH and RSA are examined in sem2,
% lectures 1 and 5).
\begin{lecture}{1}{Введение}

\section{Что такое компьютерная алгебра}

Строгого общепринятого определения у компьютерной алгебры нет. Проще всего объяснить её
предмет, сравнив с соседней дисциплиной.

\emph{Численные методы} (вычислительная математика, numerical analysis) решают задачи
\emph{в числах}: результат "--- приближённое значение, представленное числом с плавающей
точкой, а качество метода измеряется погрешностью и скоростью сходимости. \emph{Системы
компьютерной алгебры} (СКА) работают с математическими объектами \emph{точно}: с целыми и
рациональными числами произвольной длины, вычетами по модулю, многочленами, алгебраическими
числами, элементами конечных полей и групп. Ответ СКА "--- это \emph{формула} или точное
значение, а не приближение.

\begin{example}\label{l1:symbolic}
\begin{enumerate}
  \item Уравнение $x^2 - 2 = 0$ численный метод решает как $x \approx 1{,}41421356$, а СКА
    отвечает $x = \pm\sqrt 2$ и умеет дальше точно вычислять с этим объектом:
    $(1 + \sqrt 2)^2 = 3 + 2\sqrt 2$.
  \item В арифметике с плавающей точкой двойной точности
    \[0{,}1 + 0{,}2 = 0{,}30000000000000004 \ne 0{,}3,\]
    потому что десятичные дроби $0{,}1$ и $0{,}2$ не
    представимы конечной двоичной дробью. В СКА $\frac1{10} + \frac2{10} = \frac3{10}$ точно.
  \item Многочлен $x^4 - 1$ над $\QQ$ раскладывается на неприводимые множители
    $(x - 1)(x + 1)(x^2 + 1)$, а число Ферма $F_6 = 2^{64} + 1$ "--- на простые множители
    $274177 \cdot 67280421310721$ (Ф.~Ландри, 1880). Это \emph{точные} ответы, которые
    приближёнными вычислениями не получить.
\end{enumerate}
\end{example}

Примеры СКА: SageMath (свободная система на основе Python, которой мы будем пользоваться),
PARI/GP, Maxima, Maple, Mathematica, Magma. В основании любой такой системы лежит библиотека
\emph{арифметики больших чисел} (например, GMP), и алгоритмы этой арифметики "--- первая
большая тема курса.

В этом курсе компьютерная алгебра изучается с одной конкретной целью "--- как
\emph{математическая основа криптографии с открытым ключом}. Криптографии нужны именно точные
вычисления: шифртекст, искажённый в одном бите, расшифровывается в мусор, поэтому никакие
«погрешности» недопустимы. Нам понадобятся:
\begin{itemize}
  \item быстрые алгоритмы точной арифметики (умножение, деление, возведение в степень по модулю);
  \item алгебраические структуры "--- кольца вычетов $\Zn{n}$, конечные поля $\GF{p}$, группы
    точек эллиптических кривых, решётки, "--- в которых эти вычисления происходят;
  \item \emph{вычислительно сложные задачи} в этих структурах и лучшие известные алгоритмы их
    решения: именно сложность задач обеспечивает стойкость криптосистем.
\end{itemize}

\checkyourself{Приведите пример задачи, ответ на которую численный метод даёт лишь
приближённо, а система компьютерной алгебры "--- точно. Почему для криптографии нужен именно
точный ответ?}

\section{Криптографическая постановка задачи}

\subsection{Алиса, Боб и Ева}

Традиционные персонажи криптографических протоколов: \emph{Алиса} хочет передать сообщение
\emph{Бобу} по \emph{открытому каналу связи} (интернет, радиоэфир), а \emph{Ева}
(eavesdropper "--- подслушивающая) может читать всё, что передаётся по каналу
(рис.~\ref{l1:fig-channel}). Ева в нашей модели \emph{пассивна}: она слушает, но не изменяет
сообщения (активный противник, способный подменять сообщения, обычно зовётся Мэллори).

\begin{figure}[ht]
\centering
\begin{tikzpicture}[x=1cm,y=1cm]
  \node (A) at (0,0) {\personicon{caBlue}{Алиса}};
  \node (B) at (9,0) {\personicon{caRed}{Боб}};
  \draw[very thick,caGray,{Stealth[length=3mm]}-{Stealth[length=3mm]}] (1,0.3) -- (8,0.3)
    node[midway,above,font=\small] {открытый канал: $C = E_K(M)$};
  \node (E) at (4.5,-1.9) {\personicon{caGray}{Ева}};
  \draw[dashed,caGray] (4.5,0.3) -- (4.5,-1.2);
  \node[font=\small,align=left,anchor=north] at (0,-1.0) {$K$, $M$\\ $C \gets E_K(M)$};
  \node[font=\small,align=left,anchor=north] at (9,-1.0) {$K$\\ $M = D_K(C)$};
  \node[font=\small,anchor=north] at (4.5,-2.75) {видит только $C$};
\end{tikzpicture}
\caption{Симметричное шифрование: общий ключ $K$ есть у Алисы и Боба, но не у Евы}
\label{l1:fig-channel}
\end{figure}

\subsection{Симметричное шифрование}

\begin{definition}\label{l1:cipher}
\emph{Симметричный шифр} задаётся множествами ключей $\mathcal K$, открытых текстов
$\mathcal M$ и шифртекстов $\mathcal C$ и двумя отображениями
\[
  E\colon \mathcal K \times \mathcal M \to \mathcal C, \qquad
  D\colon \mathcal K \times \mathcal C \to \mathcal M,
\]
которые записывают как $E_K(M)$ (шифрование) и $D_K(C)$ (расшифрование), причём
\[
  D_K(E_K(M)) = M \quad \text{для всех } K \in \mathcal K,\ M \in \mathcal M.
\]
\end{definition}

Протокол на рис.~\ref{l1:fig-channel}: Алиса и Боб заранее владеют общим ключом $K$; Алиса
вычисляет $C = E_K(M)$, например для $M = $ «привет, что делаешь?», и отправляет $C$; Боб
вычисляет $\overline M = D_K(C)$ и по определению шифра получает $\overline M = M$.

Требования к шифру:
\begin{enumerate}
  \item $E_K$ и $D_K$ \emph{быстро} вычисляются при известном $K$;
  \item зная только шифртексты $C_1, C_2, \dots$ (и даже какие-то пары $(M_i, C_i)$), Ева не
    может \emph{за разумное время} найти $K$ или открытый текст нового сообщения.
\end{enumerate}
По \emph{принципу Керкгоффса} стойкость не должна зависеть от секретности самого алгоритма:
Ева знает $E$ и $D$, секретом является только ключ. В частности, ключей должно быть так много,
чтобы их перебор был невозможен: сегодня $\abs{\mathcal K} \ge 2^{128}$ (шифр AES).

\begin{example}
Простейший стойкий шифр "--- \emph{одноразовый блокнот}: $\mathcal K = \mathcal M =
\mathcal C = \bits{n}$, $E_K(M) = M \oplus K$, $D_K(C) = C \oplus K$ (побитовое
«исключающее или»). Корректность: $(M \oplus K) \oplus K = M$. Если $K$ выбран случайно
равномерно и используется \emph{один раз}, шифртекст не несёт никакой информации об $M$, но
ключ приходится передавать такой же длины, как все сообщения. Генераторы псевдослучайных
чисел (лекция~9) позволяют «растянуть» короткий ключ в длинную гамму.
\end{example}

\subsection{Проблема распределения ключей}

Все симметричные шифры предполагают, что \emph{общий ключ $K$ уже есть} у Алисы и Боба. Но как
его получить, если встретиться лично нельзя, а канал прослушивается? Кроме того, в сети из $n$
участников, где каждый должен уметь секретно общаться с каждым, требуется $\binom n2 =
n(n - 1)/2$ ключей. До 1976 г. эта проблема решалась организационно (курьеры, центры
распределения ключей). Решение, найденное У.~Диффи и М.~Хеллманом~\cite{diffie-hellman},
положило начало \emph{криптографии с открытым ключом}: два участника могут \emph{согласовать}
общий секретный ключ, обмениваясь сообщениями только по открытому каналу.

\checkyourself{Почему одноразовый блокнот нельзя использовать дважды с одним и тем же ключом?
Что узнает Ева из $C_1 \oplus C_2$?}

\section{Протокол Диффи--Хеллмана}\label{l1:dh}

Пусть $p$ "--- большое простое число. Мультипликативная группа $\Znx{p} = \{1, 2, \dots,
p - 1\}$ поля $\GF{p}$ циклическая (лекция~7): существует \emph{первообразный корень} $g$,
степени которого $g^0, g^1, \dots, g^{p-2}$ пробегают все ненулевые вычеты по модулю $p$.

\begin{algorithm}[протокол Диффи"--~Хеллмана]\label{l1:dh-alg}
\begin{algorithmic}[1]
\Require общедоступные параметры: простое $p$ и элемент $g \in \Znx{p}$ большого порядка
  (например, первообразный корень).
\Ensure общий секретный ключ $K$ Алисы и Боба.
\State Алиса выбирает случайное $a \in \{1, \dots, p - 2\}$ и отправляет Бобу
  $A = g^a \bmod p$
\State Боб выбирает случайное $b \in \{1, \dots, p - 2\}$ и отправляет Алисе
  $B = g^b \bmod p$
\State Алиса вычисляет $K = B^a \bmod p$, Боб вычисляет $K = A^b \bmod p$
\end{algorithmic}
\end{algorithm}

\begin{figure}[ht]
\centering
\begin{tikzpicture}[x=1cm,y=1cm]
  \node at (4.5,1.6) {общие параметры: $p$, $g$};
  \node (A) at (0,0) {\personicon{caBlue}{Алиса}};
  \node (B) at (9,0) {\personicon{caRed}{Боб}};
  \draw[-{Stealth[length=3mm]},thick] (1,0.6) -- (8,0.6) node[midway,above] {$A = g^a$};
  \draw[{Stealth[length=3mm]}-,thick] (1,0.0) -- (8,0.0) node[midway,below] {$B = g^b$};
  \node (E) at (4.5,-2.0) {\personicon{caGray}{Ева}};
  \draw[dashed,caGray] (4.5,-0.75) -- (4.5,-1.3);
  \node[font=\small,align=left,anchor=north] at (0,-1.0)
    {секрет: $a$\\ знает: $a, A, B$\\ $K \gets B^a \bmod p$};
  \node[font=\small,align=left,anchor=north] at (9,-1.0)
    {секрет: $b$\\ знает: $b, A, B$\\ $K \gets A^b \bmod p$};
  \node[font=\small,anchor=north] at (4.5,-2.9) {знает: $p, g, A, B$};
\end{tikzpicture}
\caption{Протокол Диффи"--~Хеллмана; все вычисления "--- по модулю $p$}
\label{l1:fig-dh}
\end{figure}

\begin{proposition}[корректность]\label{l1:dh-correct}
Алиса и Боб получают один и тот же ключ $K = g^{ab} \bmod p$.
\end{proposition}

\begin{proof}
Возведение в степень в группе $\Znx{p}$ подчиняется правилу $(x^u)^v = x^{uv}$, а
умножение целых чисел коммутативно, поэтому по модулю $p$
\[
  B^a \equiv (g^b)^a = g^{ab} = g^{ba} = (g^a)^b \equiv A^b. \qedhere
\]
\end{proof}

\begin{example}[\cite{hps}, пример~2.7]\label{l1:dh-example}
Пусть $p = 941$, $g = 627$ (первообразный корень: $\ord g = 940$). Алиса выбирает $a = 347$ и
отправляет $A = 627^{347} \bmod 941 = 390$; Боб выбирает $b = 781$ и отправляет
$B = 627^{781} \bmod 941 = 691$. Алиса вычисляет $691^{347} \bmod 941 = 470$, Боб "---
$390^{781} \bmod 941 = 470$. Общий ключ $K = 470$.
\end{example}

\subsection{Что знает Ева}

Ева видит $p$, $g$, $A = g^a$ и $B = g^b$. Чтобы узнать $K$, ей нужно решить следующую задачу.

\begin{definition}\label{l1:dhp}
Пусть $G$ "--- циклическая группа с образующим $g$.
\begin{enumerate}
  \item \emph{Задача дискретного логарифмирования} (DLP): по $g$ и $h \in G$ найти такое
    $x$, что $g^x = h$. Число $x$ (определённое по модулю $\ord g$) называется
    \emph{дискретным логарифмом} $h$ по основанию $g$ и обозначается $\log_g h$.
  \item \emph{Задача Диффи"--~Хеллмана} (DHP): по $g$, $g^a$, $g^b$ найти $g^{ab}$.
\end{enumerate}
\end{definition}

\begin{proposition}\label{l1:dhp-dlp}
Если Ева умеет решать DLP в группе $G$, то она умеет решать и DHP в $G$; дополнительные
затраты "--- одно возведение в степень.
\end{proposition}

\begin{proof}
Ева находит $a = \log_g A$ и вычисляет $K = B^a = g^{ab}$, как это делает Алиса.
\end{proof}

Обратное утверждение "--- можно ли решить DLP, умея решать DHP, "--- в общем случае
\emph{неизвестно}. Поэтому формально стойкость протокола опирается на сложность DHP, а
практически "--- на то, что для DHP в хорошо выбранных группах не известно способа лучше, чем
найти дискретный логарифм. Алгоритмы дискретного логарифмирования "--- полный перебор, алгоритм
Шенкса, алгоритм Полига"--~Хеллмана, $\rho$-метод Полларда, исчисление индексов "--- изучаются
в лекциях~1--4 второго семестра. Из алгоритма Полига"--~Хеллмана будет видно, что порядок $g$
обязан иметь большой простой делитель, а из исчисления индексов "--- что $p$ должно быть очень
большим (2048 бит и более).

\begin{remark}
Против \emph{активного} противника протокол в такой форме беззащитен: Мэллори, стоящий
посередине канала, согласует один ключ с Алисой, другой "--- с Бобом и перешифровывает все
сообщения (атака «человек посередине»). На практике сообщения $A$, $B$ подписываются
(электронная подпись), но это выходит за рамки курса.
\end{remark}

\checkyourself{Почему в протоколе Диффи"--~Хеллмана нельзя брать $g = 1$ или $g = p - 1$?
Что увидит Ева в этих случаях?}

\checkyourself{Ева перехватила $A$ и $B$ и знает $a$. Может ли она вычислить $K$? А если
она знает только произведение $ab$?}

\section{Односторонние функции}

Протокол Диффи"--~Хеллмана работает потому, что функция $x \mapsto g^x \bmod p$
\emph{вычисляется быстро} (см. ниже алгоритм~\ref{l1:binpow}), а обратная к ней "---
дискретный логарифм "--- \emph{вычисляется трудно}. Это свойство заслуживает отдельного
названия.

\begin{definition}[неформальное]\label{l1:owf}
Функция $f\colon X \to Y$ называется \emph{односторонней}, если
\begin{enumerate}
  \item значение $y = f(x)$ \emph{легко} вычисляется для любого $x \in X$;
  \item для случайно выбранного $x$ и $y = f(x)$ найти \emph{какой-нибудь} прообраз
    $x' \in f^{-1}(y)$ "--- \emph{вычислительно сложная задача}.
\end{enumerate}
\end{definition}

Слова «легко» и «трудно» уточняются в теории сложности вычислений: «легко» "--- за время,
полиномиальное от длины входа (в битах), «трудно» "--- любой полиномиальный (вероятностный)
алгоритм находит прообраз лишь с пренебрежимо малой вероятностью. Точные определения см.
в~\cite[разд.~1.3]{hac} и~\cite[гл.~8]{hps}; нам будет достаточно неформального понимания,
подкреплённого оценками сложности \emph{лучших известных алгоритмов} (разд.~\ref{l1:complexity}).

\begin{remark}
\emph{Существование односторонних функций не доказано.} Более того, из их существования
следовало бы $\mathrm P \ne \mathrm{NP}$: проверка «$x$ "--- прообраз $y$» выполняется за
полиномиальное время, поэтому при $\mathrm P = \mathrm{NP}$ прообраз можно было бы найти
полиномиальным алгоритмом (восстанавливая $x$ по одному биту). Вся криптография с открытым
ключом опирается на \emph{гипотезы} о сложности конкретных задач.
\end{remark}

\subsection{Кандидаты в односторонние функции}

\begin{enumerate}
  \item \textbf{Модульное возведение в степень}: $x \mapsto g^x \bmod p$. Обращение "---
    задача дискретного логарифмирования. Аналог в группе точек эллиптической кривой: $k
    \mapsto kP$, обращение "--- задача ECDLP (лекции~8--9 второго семестра).
  \item \textbf{Умножение}: $(p, q) \mapsto N = pq$ для простых $p$, $q$ одинаковой длины.
    Обращение "--- \emph{факторизация} $N$. Перебор делителей до $\sqrt N$ требует порядка
    $\sqrt N$ делений; лучшие известные алгоритмы субэкспоненциальны (лекции~5, 6 и~10
    второго семестра).
  \item \textbf{Возведение в квадрат по составному модулю}: $x \mapsto x^2 \bmod N$, $N = pq$.
    Обращение "--- извлечение квадратного корня по модулю $N$; эта задача равносильна
    факторизации $N$ (см. утверждение~\ref{l1:sqrt-factor}).
  \item \textbf{Задачи на решётках}: по базису решётки $\Lat \subset \RR^n$ найти её
    кратчайший ненулевой вектор (SVP) или ближайший к заданной точке вектор решётки (CVP).
    Функция «точка решётки + малый шум $\mapsto$ результат» обращается решением CVP.
    Эти задачи считаются трудными даже для квантового компьютера, на котором DLP и
    факторизация решаются полиномиально (алгоритм Шора); см. лекции~11--12 второго семестра.
\end{enumerate}

\begin{lemma}\label{l1:sqrt-gcd}
Пусть $N$ "--- натуральное, $\congr{x^2}{y^2}{N}$, но $x \not\equiv \pm y \Mod N$. Тогда
$d = \gcd(x - y, N)$ "--- нетривиальный делитель $N$: $1 < d < N$.
\end{lemma}

\begin{proof}
По условию $N \mid x^2 - y^2 = (x - y)(x + y)$. Если $d = N$, то $N \mid x - y$, т.\,е.
$\congr{x}{y}{N}$ "--- противоречие. Если $d = 1$, то $N$ взаимно просто с $x - y$ и потому
делит $x + y$, т.\,е. $\congr{x}{-y}{N}$ "--- снова противоречие.
\end{proof}

\begin{proposition}\label{l1:sqrt-factor}
Пусть $N = pq$, где $p \ne q$ "--- нечётные простые. Если есть алгоритм $\mathcal A$, который
по квадрату $c \in \Znx{N}$ возвращает какой-нибудь его квадратный корень, то $N$ можно
разложить на множители, обращаясь к $\mathcal A$ в среднем не более двух раз.
\end{proposition}

\begin{proof}
Сначала покажем, что у каждого квадрата $c = x^2 \in \Znx{N}$ ровно четыре квадратных корня.
По модулю простого $p$ сравнение $y^2 \equiv x^2$ равносильно $p \mid (y - x)(y + x)$, т.\,е.
$y \equiv \pm x \Mod p$, и это два \emph{разных} вычета, так как $p \nmid 2x$. То же по модулю
$q$. По китайской теореме об остатках (лекция~2 второго семестра) каждой из $2 \cdot 2 = 4$
комбинаций знаков соответствует ровно один вычет $y$ по модулю $N$; среди них $\pm x$ и ещё
два корня $\pm x'$ с $x' \not\equiv \pm x$.

Алгоритм: выберем случайное $x \in \Znx{N}$ (если случайно $\gcd(x, N) > 1$, делитель уже
найден), вычислим $c = x^2 \bmod N$ и $y = \mathcal A(c)$. Ответ $\mathcal A$ зависит только
от $c$, а $c$ одинаково получается из всех четырёх своих корней, которые выбирались бы
равновероятно; значит, с вероятностью $1/2$ окажется $y \not\equiv \pm x$, и по
лемме~\ref{l1:sqrt-gcd} $\gcd(x - y, N)$ "--- нетривиальный делитель. Среднее число попыток до
первого успеха при вероятности успеха $1/2$ равно $2$.
\end{proof}

\begin{example}
При $N = 77 = 7 \cdot 11$ у вычета $4$ четыре квадратных корня: $2$, $9$, $68$, $75$
(действительно, $9^2 = 81 = 77 + 4$). Зная корни $2$ и $9 \not\equiv \pm 2$, получаем
$\gcd(9 - 2, 77) = 7$ и $\gcd(9 + 2, 77) = 11$.
\end{example}

Обратно, если известны $p$ и $q$, квадратные корни по модулю $p$ и $q$ извлекаются быстро
(алгоритмы Тонелли"--~Шенкса и Чиполлы, лекция~7 второго семестра) и склеиваются по китайской
теореме об остатках. Таким образом, разложение $N$ служит \emph{секретом}, облегчающим
обращение функции $x \mapsto x^2 \bmod N$.

\subsection{Односторонние функции с секретом}

\begin{definition}[неформальное]\label{l1:tdf}
\emph{Односторонняя функция с потайным входом} (с секретом, trapdoor one-way function) "---
это односторонняя функция $f$, для которой существует дополнительная информация $t$
(\emph{секрет}, trapdoor), зная которую прообраз $f^{-1}(y)$ можно найти эффективно.
\end{definition}

Для чего нужен секрет? Он превращает одностороннюю функцию в \emph{асимметричный шифр}:
функцию $f$ (точнее, её описание "--- \emph{открытый ключ}) публикуют, и любой может
\emph{зашифровать} сообщение $m$, вычислив $c = f(m)$; обратить $f$ и \emph{расшифровать}
$c$ может только владелец секрета $t$ "--- \emph{закрытого ключа}. Проблема распределения
ключей при этом исчезает: Алисе не нужно заранее договариваться с Бобом "--- достаточно взять
его открытый ключ из справочника. Главный пример "--- RSA.

\checkyourself{Чем односторонняя функция отличается от односторонней функции с секретом?
Какой секрет у функции $x \mapsto x^2 \bmod N$? Почему функция $x \mapsto g^x \bmod p$
не годится для асимметричного шифрования так же непосредственно?}

\section{RSA как функция с секретом}\label{l1:rsa}

Криптосистема RSA (Р.~Ривест, А.~Шамир, Л.~Адлеман, 1977~\cite{rsa}) основана на сложности
факторизации и на теореме Эйлера.

\subsection{Функция Эйлера и теорема Эйлера}

\begin{definition}
\emph{Функция Эйлера} $\varphi(n)$ "--- число вычетов $x \in \{0, 1, \dots, n - 1\}$,
взаимно простых с $n$, т.\,е. порядок группы $\Znx{n}$.
\end{definition}

\begin{lemma}\label{l1:phi-pq}
Для простого $p$ имеем $\varphi(p) = p - 1$; для различных простых $p$, $q$
\[
  \varphi(pq) = (p - 1)(q - 1).
\]
\end{lemma}

\begin{proof}
Число $x \in \{0, \dots, pq - 1\}$ не взаимно просто с $pq$ тогда и только тогда, когда оно
делится на $p$ или на $q$. Кратных $p$ среди них $q$ штук ($0, p, \dots, (q-1)p$), кратных
$q$ "--- $p$ штук, а кратно обоим (т.\,е. $pq$) только число $0$. По формуле
включений"--~исключений
$\varphi(pq) = pq - q - p + 1 = (p - 1)(q - 1)$. Случай $\varphi(p) = p - 1$ очевиден.
\end{proof}

\begin{theorem}[Эйлер]\label{l1:euler}
Если $\gcd(a, n) = 1$, то $\congr{a^{\varphi(n)}}{1}{n}$.
\end{theorem}

\begin{proof}
Пусть $r_1, \dots, r_k$, $k = \varphi(n)$, "--- все вычеты по модулю $n$, взаимно простые с $n$.
Числа $a r_1, \dots, a r_k$ также взаимно просты с $n$ и попарно несравнимы: из
$\congr{a r_i}{a r_j}{n}$ следует $n \mid a(r_i - r_j)$, а так как $\gcd(a, n) = 1$, то
$n \mid r_i - r_j$, т.\,е. $i = j$. Значит, $a r_1, \dots, a r_k$ "--- те же вычеты
$r_1, \dots, r_k$, взятые в другом порядке, и
\[
  a^k \, r_1 r_2 \cdots r_k \equiv (a r_1)(a r_2) \cdots (a r_k) \equiv r_1 r_2 \cdots r_k
  \Mod n.
\]
Произведение $r_1 \cdots r_k$ взаимно просто с $n$, поэтому на него можно сократить:
$\congr{a^k}{1}{n}$.
\end{proof}

При $n = p$ простом получаем \emph{малую теорему Ферма}: $\congr{a^{p-1}}{1}{p}$ при
$p \nmid a$; умножив на $a$, имеем $\congr{a^p}{a}{p}$ уже для \emph{всех} целых $a$.

\subsection{Схема RSA}

\begin{algorithm}[генерация ключей RSA]\label{l1:rsa-keygen}
\begin{algorithmic}[1]
\Require длина ключа (например, $2048$ бит).
\Ensure открытый ключ $(e, N)$, закрытый ключ $(d, p, q)$.
\State выбрать случайные простые $p \ne q$ половинной длины \Comment{тесты простоты, лекция~8}
\State $N \gets pq$, \quad $\varphi(N) \gets (p - 1)(q - 1)$
\State выбрать $e$ с $\gcd(e, \varphi(N)) = 1$ \Comment{на практике часто $e = 65537$}
\State $d \gets e^{-1} \bmod \varphi(N)$ \Comment{расширенный алгоритм Евклида, лекция~3}
\State \Return $(e, N)$, $(d, p, q)$
\end{algorithmic}
\end{algorithm}

\emph{Шифрование} сообщения $m \in \{0, 1, \dots, N - 1\}$ открытым ключом $(e, N)$:
$c = m^e \bmod N$. \emph{Расшифрование} закрытым ключом: $m = c^d \bmod N$. Таким образом,
функция с секретом "--- это $f(m) = m^e \bmod N$, а секрет "--- число $d$ или, что то же самое,
разложение $N = pq$ (зная $p$ и $q$, $d$ находится алгоритмом Евклида).

\begin{theorem}[корректность RSA]\label{l1:rsa-correct}
Пусть $N = pq$, $p \ne q$ простые, $\congr{ed}{1}{\varphi(N)}$. Тогда
$\congr{m^{ed}}{m}{N}$ для \emph{всех} целых $m$.
\end{theorem}

\begin{proof}
Запишем $ed = 1 + k(p - 1)(q - 1)$, $k \ge 0$. Докажем, что $\congr{m^{ed}}{m}{p}$.
Если $p \mid m$, обе части сравнимы с нулём. Если $p \nmid m$, то по теореме~\ref{l1:euler}
(малой теореме Ферма) $\congr{m^{p-1}}{1}{p}$ и
\[
  m^{ed} = m \cdot \bigl(m^{p-1}\bigr)^{k(q-1)} \equiv m \cdot 1 = m \Mod p.
\]
Точно так же $\congr{m^{ed}}{m}{q}$. Итак, $p \mid m^{ed} - m$ и $q \mid m^{ed} - m$, а так
как $p$ и $q$ "--- различные простые, то и $pq = N \mid m^{ed} - m$.
\end{proof}

\begin{remark}
Если $\gcd(m, N) = 1$, утверждение сразу следует из теоремы Эйлера:
$m^{ed} = m \cdot (m^{\varphi(N)})^k \equiv m$. Разбор по простым делителям нужен лишь для
«неудачных» $m$, кратных $p$ или $q$ (их доля $(p + q - 1)/N$ ничтожна, но корректность
должна выполняться для всех сообщений).
\end{remark}

\begin{example}\label{l1:rsa-example}
Пусть $p = 47$, $q = 59$. Тогда $N = 2773$, $\varphi(N) = 46 \cdot 58 = 2668$. Возьмём $e = 17$
($\gcd(17, 2668) = 1$); тогда $d = 17^{-1} \bmod 2668 = 157$: действительно,
$17 \cdot 157 = 2669 = 2668 + 1$. Открытый ключ $(17, 2773)$, закрытый "--- $(157, 47, 59)$.
\begin{itemize}
  \item Шифрование $m = 920$: $c = 920^{17} \bmod 2773 = 948$; расшифрование:
    $948^{157} \bmod 2773 = 920$.
  \item «Неудачное» сообщение $m = 94 = 2 \cdot 47$: $c = 94^{17} \bmod 2773 = 1833$, и
    $1833^{157} \bmod 2773 = 94$ "--- как и утверждает теорема~\ref{l1:rsa-correct}.
\end{itemize}
\end{example}

\subsection{Обмен сообщениями и стойкость}

Если у каждого участника своя пара ключей "--- у Алисы $(e_A, N_A)$ и $(d_A, p_A, q_A)$, у
Боба $(e_B, N_B)$ и $(d_B, p_B, q_B)$, "--- то Алиса шифрует сообщение $m_1$ \emph{открытым
ключом Боба}: $c_1 = m_1^{e_B} \bmod N_B$, и только Боб может вычислить
$m_1 = c_1^{d_B} \bmod N_B$. Ответ $m_2$ Боб шифрует открытым ключом Алисы:
$c_2 = m_2^{e_A} \bmod N_A$, и Алиса расшифровывает: $m_2 = c_2^{d_A} \bmod N_A$.

Ева знает $N$, $e$ и $c$. Если она разложит $N$, то вычислит $\varphi(N)$ и $d$, т.\,е.
\emph{RSA не стойче факторизации}. Верно и следующее.

\begin{proposition}\label{l1:phi-factor}
Знание $\varphi(N)$ для $N = pq$ равносильно знанию разложения $N$.
\end{proposition}

\begin{proof}
Если известны $p$ и $q$, то $\varphi(N) = (p - 1)(q - 1)$. Обратно, $\varphi(N) = N - (p + q) +
1$, откуда $s = p + q = N - \varphi(N) + 1$, а $p$, $q$ "--- корни квадратного уравнения
$z^2 - sz + N = 0$: $p, q = \bigl(s \pm \sqrt{s^2 - 4N}\bigr)/2$.
\end{proof}

Для примера~\ref{l1:rsa-example}: $s = 2773 - 2668 + 1 = 106$, $s^2 - 4N = 11236 - 11092 = 144$,
$p, q = (106 \pm 12)/2 = 59, 47$. Можно также показать (без доказательства;
см.~\cite[разд.~3.2]{hps} и упражнения к нему), что и знание $d$ позволяет эффективно
(вероятностным алгоритмом) разложить $N$. Открытым остаётся вопрос, можно ли извлекать корни
$e$-й степени по модулю $N$, \emph{не} раскладывая $N$.

\begin{remark}
Описанная схема "--- «учебная» (textbook) RSA. Она детерминирована (одинаковые сообщения
дают одинаковые шифртексты) и мультипликативна ($f(m_1 m_2) = f(m_1) f(m_2)$), поэтому на
практике перед шифрованием сообщение дополняют случайными данными по стандарту (OAEP).
Атаки на RSA и факторизация "--- тема лекций~5, 6 и~10 второго семестра.
\end{remark}

\checkyourself{В примере~\ref{l1:rsa-example} проверьте, что $e = 23$ не годится в качестве
открытой экспоненты. Почему $e$ должно быть взаимно просто с $\varphi(N)$?}

\checkyourself{Почему в доказательстве теоремы~\ref{l1:rsa-correct} существенно, что $p \ne q$?
Верно ли утверждение теоремы для $N = p^2$ и $\varphi(N) = p(p - 1)$?}

\section{Сложность: почему нужны большие числа}\label{l1:complexity}

\subsection{Размер входа}

Сложность алгоритма измеряется как функция \emph{длины входа в битах}. Натуральное число $N$
занимает
\[
  n = \floor{\log_2 N} + 1 \approx \log_2 N
\]
битов (лекция~2), а \emph{не} $N$ единиц памяти. Это принципиально: алгоритм, выполняющий
$\sqrt N$ операций, кажется «быстрым» как функция от $N$, но $\sqrt N = 2^{n/2}$
экспоненциально по длине входа. Например, модуль RSA длины $n = 2048$ бит "--- это число из
$617$ десятичных цифр, и перебор делителей до $\sqrt N \approx 2^{1024}$ невозможен.

\begin{definition}\label{l1:classes}
Пусть алгоритм обрабатывает входы длины $n$ за время $T(n)$. Говорят, что алгоритм
\begin{itemize}
  \item \emph{полиномиальный}, если $T(n) = \bigO(n^k)$ для некоторой константы $k$;
  \item \emph{экспоненциальный}, если $T(n) = 2^{\Theta(n)}$ (например, $T = \bigO(\sqrt N) =
    \bigO(2^{n/2})$);
  \item \emph{субэкспоненциальный}, если $T(n) = 2^{o(n)}$, но $T(n)$ растёт быстрее любого
    полинома от $n$.
\end{itemize}
\end{definition}

Полиномиальные алгоритмы считаются «быстрыми», экспоненциальные "--- непрактичными для больших
входов. Субэкспоненциальные занимают промежуточное положение и для криптографии важнее всего:
именно такие алгоритмы известны для факторизации и для DLP в $\Znx{p}$. Их сложность удобно
записывать с помощью \emph{L-нотации}:
\[
  L_N[\alpha, c] = \exp\Bigl( (c + o(1)) \, (\ln N)^{\alpha} \, (\ln \ln N)^{1 - \alpha} \Bigr),
  \qquad 0 \le \alpha \le 1,\ c > 0.
\]
Параметр $\alpha$ интерполирует между полиномиальной и экспоненциальной сложностью:
\[
  L_N[0, c] = (\ln N)^{c + o(1)} \ \text{(полином от $n$)}, \qquad
  L_N[1, c] = N^{c + o(1)} = 2^{(c + o(1)) n} \ \text{(экспонента)}.
\]
Действительно, при $\alpha = 0$ показатель равен $(c + o(1)) \ln \ln N$, а
$e^{c \ln \ln N} = (\ln N)^c$; при $\alpha = 1$ показатель равен $(c + o(1)) \ln N$.

\subsection{Лёгкое направление: быстрое возведение в степень}

Для протоколов важно, что «законные» операции полиномиальны. Возведение в степень $g^a \bmod p$
наивно требует $a - 1$ умножений, т.\,е. экспоненциального числа операций, но его можно выполнить
за $\bigO(\log a)$ умножений.

\begin{algorithm}[бинарное возведение в степень справа налево]\label{l1:binpow}
\begin{algorithmic}[1]
\Require $g \in \Zn{m}$, $a \ge 0$.
\Ensure $g^a \bmod m$.
\State $r \gets 1$, \quad $s \gets g$
\While{$a > 0$}
  \If{$a$ нечётно} \State $r \gets r s \bmod m$ \EndIf
  \State $s \gets s^2 \bmod m$, \quad $a \gets \floor{a/2}$
\EndWhile
\State \Return $r$
\end{algorithmic}
\end{algorithm}

\begin{proposition}\label{l1:binpow-cost}
Алгоритм~\ref{l1:binpow} вычисляет $g^a \bmod m$ не более чем за $2\ell$ умножений по
модулю $m$, где $\ell = \floor{\log_2 a} + 1$ "--- битовая длина $a$. Для $n$-битных $m$ и $a$
это $\bigO(n^3)$ битовых операций при школьном умножении.
\end{proposition}

\begin{proof}
Пусть $a_0$ "--- исходное значение показателя, а $a$, $r$, $s$ "--- текущие. Инвариант
цикла: $r \cdot s^{a} \equiv g^{a_0} \Mod m$. Вначале $r = 1$, $s = g$, $a = a_0$. Если
$a = 2t$, то после шага $r s^a = r (s^2)^t$ "--- инвариант сохраняется; если $a = 2t + 1$, то
$r s^a = (rs)(s^2)^t$ "--- тоже. Цикл заканчивается при $a = 0$, и тогда $r \equiv g^{a_0}$.
Каждая итерация уменьшает битовую длину $a$ на единицу, поэтому итераций ровно $\ell$, и в каждой
не более двух умножений. Умножение $n$-битных чисел с приведением по модулю стоит $\bigO(n^2)$
битовых операций (лекции~2--3), откуда $\bigO(n \cdot n^2)$.
\end{proof}

Подробно методы возведения в степень разбираются в лекции~6, быстрое умножение "--- в лекциях~2
и~4, модулярная арифметика (метод Монтгомери) "--- в лекции~5.

\subsection{Трудное направление и выбор длины ключа}

В таблице~\ref{l1:tab-costs} собраны лучшие известные алгоритмы для задач, на которых
основаны DH и RSA (здесь $N$ "--- модуль RSA, $p$ "--- модуль DH, $q$ "--- порядок группы точек
эллиптической кривой).

\begin{table}[!htb]
\centering\small
\caption{Сложность лучших известных (неквантовых) алгоритмов}\label{l1:tab-costs}
\begin{tabular}{p{0.37\linewidth}lp{0.2\linewidth}}
\hline
Задача и алгоритм & Сложность & Где в курсе \\
\hline
факторизация: перебор делителей & $\bigO(\sqrt N) = L_N[1, \tfrac12]$ & сем.~2, лекция~5 \\
DLP в любой группе: Шенкс, $\rho$-метод & $\bigO(\sqrt q)$ & сем.~2, лекции~1, 3 \\
DLP в $\Znx{p}$: исчисление индексов & $L_p[\tfrac12, c]$ & сем.~2, лекция~4 \\
факторизация: квадратичное решето & $L_N[\tfrac12, 1]$ & сем.~2, лекция~6 \\
факторизация: метод Ленстры & $L_p[\tfrac12, \sqrt2]$, $p \mid N$ & сем.~2, лекция~10 \\
решето числового поля (факторизация, DLP) & $L_N[\tfrac13, \sqrt[3]{64/9}]$ & обзорно \\
ECDLP на «хорошей» кривой & $\bigO(\sqrt q)$ (лучше не известно) & сем.~2, лекция~9 \\
\hline
\end{tabular}
\end{table}

Чтобы почувствовать разницу, подставим $N = 2^n$ и \emph{отбросим} $o(1)$ в показателе (это
даёт лишь порядок величины, но наглядно); в таблице~\ref{l1:tab-bits} указан $\log_2$ числа
операций.

\begin{table}[!htb]
\centering\small
\caption{$\log_2$ числа операций при $N = 2^n$ (без учёта $o(1)$)}\label{l1:tab-bits}
\begin{tabular}{rrrr}
\hline
$n$ & $\sqrt N$ & $L_N[\tfrac12, 1]$ & $L_N[\tfrac13, 1{,}923]$ \\
\hline
512 & 256 & 66 & 64 \\
1024 & 512 & 98 & 87 \\
2048 & 1024 & 146 & 117 \\
3072 & 1536 & 184 & 139 \\
\hline
\end{tabular}
\end{table}

Выводы:
\begin{itemize}
  \item стойкость достигается \emph{большими числами}: «запас» $2^{80}$--$2^{128}$ операций
    у атакующего требует модулей RSA и DH длиной $2048$--$3072$ бит (512-битный модуль RSA был
    разложен ещё в 1999 г.);
  \item для эллиптических кривых, где известны лишь экспоненциальные $\bigO(\sqrt q)$-алгоритмы,
    достаточно $q \approx 2^{256}$ "--- поэтому ключи там в~10 раз короче;
  \item при этом законные операции "--- $\bigO(n^3)$ или быстрее "--- остаются практичными:
    именно ради них нужна эффективная арифметика больших чисел первого семестра.
\end{itemize}

На практике качественная и количественная оценка \emph{сложности задачи} определяется
сложностью \emph{лучшего известного алгоритма} её решения. Открытие нового алгоритма может
обесценить криптосистему, поэтому изучение методов решения сложных задач "--- такая же
важная часть курса, как и быстрые алгоритмы.

\checkyourself{Почему алгоритм, делающий $\bigO(N)$ операций над числом $N$, не считается
полиномиальным? Проверьте, что $L_N[0, c]$ "--- полином от $\log N$, а $L_N[1, c]$ "---
степень $N$.}

\checkyourself{Сколько умножений по модулю выполнит алгоритм~\ref{l1:binpow} при $a = 2^{k}$
и при $a = 2^{k} - 1$? Оцените число умножений при вычислении $627^{781} \bmod 941$ из
примера~\ref{l1:dh-example}.}

\section{Программа курса}

Курс посвящён \emph{алгебраическим структурам} и \emph{сложным задачам} в них, используемым
для построения криптографических примитивов. Мы изучим быстрые алгоритмы вычислений, сами
вычислительно сложные задачи и методы их решения с оценками сложности.

\begin{center}\small
\begin{tabular}{lp{0.5\linewidth}p{0.36\linewidth}}
\hline
& Тема & Зачем нужна \\
\hline
\multicolumn{3}{l}{\emph{Семестр 1. Арифметика больших чисел и элементы теории чисел (зачёт)}} \\
Л2 & Арифметика больших чисел; Карацуба, Тоом"--~Кук & все «законные» вычисления \\
Л3 & Деление с остатком и алгоритм Евклида & ключ $d$ в RSA, обратные элементы \\
Л4 & Свёртки и умножение с помощью БПФ & умножение очень больших чисел \\
Л5 & Модулярная арифметика: Монтгомери, Ферма, Мерсенн & быстрое умножение по модулю \\
Л6 & Быстрое возведение в степень & DH, RSA, тесты простоты \\
Л7 & Порядок элемента и первообразные корни & выбор $g$ в DH \\
Л8 & Простые числа и тесты простоты & генерация $p$ (DH) и $p, q$ (RSA) \\
Л9 & Генераторы псевдослучайных чисел & секреты $a$, $b$, ключи, гамма \\
\hline
\multicolumn{3}{l}{\emph{Семестр 2. Сложные задачи, эллиптические кривые, решётки (экзамен)}} \\
Л1 & Задача дискретного логарифмирования. Алгоритм Шенкса & атака на DH \\
Л2 & КТО и алгоритм Полига"--~Хеллмана & выбор порядка $g$ \\
Л3 & Алгоритмы поиска коллизий. $\rho$-алгоритм Полларда & универсальная атака на DLP \\
Л4 & Метод исчисления индексов & почему $p$ так велико \\
Л5 & Факторизация, RSA и алгоритм $p - 1$ Полларда & атаки на RSA, выбор $p, q$ \\
Л6 & Разность квадратов и квадратичное решето & субэкспоненц. факторизация \\
Л7 & Квадратичные сравнения. Методы Чиполлы и Гензеля & функция $x \mapsto x^2 \bmod N$ \\
Л8 & Эллиптические кривые & группа для ECDH \\
Л9 & Возведение точки в степень и задача ECDLP & короткие ключи \\
Л10 & Алгоритм Ленстры (факторизация кривыми) & поиск малых делителей \\
Л11 & Решётки. Алгоритмы Гаусса и Бабаи & SVP, CVP, функции на решётках \\
Л12 & Алгоритм LLL & криптоанализ, постквантовые схемы \\
\hline
\end{tabular}
\end{center}

\section{Литература к курсу}

\begin{itemize}
  \item Дж.~Хоффштейн, Дж.~Пайфер, Дж.~Сильверман~\cite{hps} "--- \textbf{основной учебник}
    курса: аккуратное и доступное введение в математическую криптографию (DLP, RSA, эллиптические
    кривые, решётки) с большим количеством упражнений. К этой лекции "--- гл.~1--3.
  \item Р.~Крэндалл, К.~Померанс~\cite{crandall-pomerance} "--- исчерпывающая информация о
    простых числах, тестах простоты, факторизации и быстрой арифметике (гл.~9).
  \item Д.~Кнут~\cite{knuth-taocp2} "--- классика: арифметика больших чисел (гл.~4) и
    генераторы случайных чисел (гл.~3) с подробным анализом алгоритмов.
  \item О.~Н.~Василенко~\cite{vasilenko} "--- обзор теоретико-числовых алгоритмов
    криптографии на русском языке: тесты простоты, факторизация, дискретное логарифмирование.
  \item В.~М.~Чернов, А.~О.~Корепанов~\cite{chernov-korepanov} "--- теоретико-числовые
    преобразования (лекции~4--5).
  \item С.~Гэлбрейт~\cite{galbraith} "--- подробный и современный справочник по математике
    криптографии с открытым ключом (свободно доступная версия на сайте автора).
  \item «Справочник по прикладной криптографии»~\cite{hac} "--- энциклопедия алгоритмов
    (свободно доступна); удобен как справочник.
  \item К.~Паар, Я.~Пельцль~\cite{paar-pelzl} "--- инженерный взгляд на криптографию, к книге
    есть видеолекции~\cite{paar-lectures}.
\end{itemize}

\checkyourselfsection

\begin{furtherreading}
  \readitem[гл.~1, разд.~1.7; гл.~2, разд.~2.1--2.3, 2.6; гл.~3, разд.~3.1--3.2]{hps}{
    симметричные и асимметричные шифры, протокол Диффи"--~Хеллмана, задачи DLP и DHP,
    O-нотация и классы сложности, теорема Эйлера и RSA}
  \readitem[гл.~1, 6, 7; разд.~8.1]{paar-pelzl}{введение в криптографию, идея открытого ключа,
    RSA и протокол Диффи"--~Хеллмана с инженерной точки зрения}
  \readitem[гл.~1; разд.~3.1]{hac}{обзор криптографии: односторонние функции и функции с
    секретом; таблица сложных теоретико-числовых задач и сводимостей между ними}
  \readitem{diffie-hellman}{оригинальная статья: идея криптографии с открытым ключом и протокол
    согласования ключа}
  \readitem{rsa}{оригинальная статья RSA}
  \readitem{coursera-boneh}{онлайн-курс Д.~Боне: модели стойкости, протокол Диффи"--~Хеллмана,
    функции с секретом и RSA (последние недели курса)}
  \readitem{paar-lectures}{видеолекции К.~Паара: введение в криптографию с открытым ключом, RSA,
    Диффи"--~Хеллман}
\end{furtherreading}
\end{lecture}
